% ATTEMPT_STATUS: unresolved
% TOP_PROBLEM: {"problemId": "problem.hyperkahler-deformation-finiteness-in-fixed-dimension", "canonicalTitle": "Boundedness of Compact Hyperkahler Manifolds: Finiteness of Deformation Types in Each Dimension", "releaseRank": 421, "primaryDomain": "geometry_and_topology", "primaryDomainLabel": "Geometry & topology", "displayStatus": "open_with_solved_subcases", "releaseStatus": "open_with_solved_subcases"}
% !TeX program = lualatex
\documentclass[11pt]{article}
\usepackage[margin=25mm]{geometry}
\usepackage{fontspec}
\setmainfont{Latin Modern Roman}
\usepackage{amsmath,amssymb,mathtools}
\usepackage{unicode-math}
\setmathfont{Latin Modern Math}
\usepackage{xurl,hyperref}
\hypersetup{colorlinks=true,urlcolor=blue}
\setcounter{secnumdepth}{0}
\setlength{\parindent}{0pt}
\setlength{\parskip}{5pt}
\setlength{\emergencystretch}{3em}
\begin{document}
\section*{\texorpdfstring{Boundedness of Compact Hyperkahler Manifolds: Finiteness of Deformation Types in Each Dimension}{Boundedness of Compact Hyperkahler Manifolds: Finiteness of Deformation Types in Each Dimension}}\label{421}
\subsection{Definitions and mathematical statement}
An irreducible holomorphic symplectic manifold $X$ is smooth, compact, simply connected and K\"ahler, and has $H^0(X,\Omega_X^2)={\mathbb{C}}\sigma$ for an everywhere nondegenerate holomorphic two-form. Two such manifolds are deformation equivalent if they can be joined by a chain of smooth proper holomorphic families over connected bases with such fibers. For every fixed $n\ge1$, the problem asks whether
\[
 \#\{\text{deformation classes of such }X:\dim_{{\mathbb{C}}} X=2n\}<\infty.
\]
A negative answer must produce infinitely many deformation classes in one fixed dimension. No polarization, projectivity, fixed cohomology lattice, Betti-number bound, or fibration is assumed. Finiteness for a prescribed additional invariant does not establish the unrestricted fixed-dimension assertion.
\par\smallskip\textit{Formulation and concept references:} \href{https://arxiv.org/pdf/1607.03215}{[S1]}.

\subsection{Short English statement}
In each fixed even complex dimension, are there only finitely many deformation types of compact irreducible holomorphic symplectic manifolds?
\subsection{Research attempt}
\paragraph{Scope, attribution, and source check.}
Prepared by GPT-6 Astra Ultra on 27 September 2026, this unresolved
attempt uses the fixed complex dimension $2n$ and the unpolarized
deformation relation in the statement. Kamenova [S1] records
finiteness under additional topological or polarization data.
Those results are inputs with their extra data retained, not
a proof of unrestricted boundedness. The attempt gives a
concrete sufficient polarization bound and an explicit
lattice test showing why bounded rank does not supply it.
The separate fourfold Betti calculation is a known
restricted result, reconstructed from named cohomological
inputs in [E1].

\paragraph{A precise polarization route.}
Suppose that, for each fixed $n$, there is an integer $D(n)$
such that every irreducible holomorphic symplectic
$2n$-fold is deformation equivalent to a projective one
$Y$ with an ample line bundle $A$ satisfying
\begin{equation}\label{a421:degree}
                       1\leq\int_Yc_1(A)^{2n}\leq D(n).
\end{equation}
Then the desired finiteness follows from the
Koll\'ar--Matsusaka boundedness input quoted in [S1]:
in a fixed dimension there are finitely many deformation
types of polarized projective manifolds with fixed
$A^{2n}$ and $K_Y\cdot A^{2n-1}$.
Here the holomorphic symplectic form trivializes $K_Y$
by its $n$-th exterior power, so the second number is zero.
There are only finitely many integer choices for
the first number in \eqref{a421:degree}.
Forgetting polarization and taking a finite union
therefore gives finitely many unpolarized
deformation types. The assumed deformation
from $X$ to $Y$ places every original manifold
in that same finite list.

This argument distinguishes an ample line bundle
from a merely positive-square or nef class.
Boundedness theorems with birational conclusions
would need additional information relating
birational holomorphic symplectic manifolds
to the stated deformation relation. The
present sufficient hypothesis avoids hiding
that extra step.

\paragraph{The Fujiki identity identifies two separate bounds.}
Normalize the Beauville--Bogomolov form $q_X$ to be the
primitive integral form with the conventional positive
cone. The established Fujiki identity says
\[
 \int_X\alpha^{2n}=c_Xq_X(\alpha)^n,\qquad c_X>0.
\]
Thus a route to \eqref{a421:degree} would be uniform
bounds $c_X\leq C(n)$ and $0<q_Y(c_1(A))\leq Q(n)$
for a suitable ample class on a deformation:
then $A^{2n}\leq C(n)Q(n)^n$.
Neither the constant nor the availability of such
an ample class has been controlled here.

There is a useful exact mixed-intersection test.
If $q(b)=0$, substitute $\alpha=a+tb$:
\[
 \int_X(a+tb)^{2n}
   =c_X\bigl(q(a)+2tq(a,b)\bigr)^n.
\]
Comparing the coefficient of $t^n$ gives
\begin{equation}\label{a421:mixed}
 \binom{2n}{n}\int_Xa^nb^n
                       =c_X\,2^nq(a,b)^n.
\end{equation}
The factor $2^n$ uses the bilinear convention
$q(a+tb)=q(a)+2tq(a,b)+t^2q(b)$.
This identity can constrain a proposed fibration
or cycle calculation, but an integral isotropic
vector alone does not produce a fibration.
One still needs it to become an appropriate
line bundle and then to satisfy the relevant
positivity and semiampleness hypotheses.

\paragraph{Bounded rank does not bound the lattice.}
Consider the integral lattices
\[
 \Lambda_D=\bigl(\mathbb Z^4,
                   \operatorname{diag}(1,1,1,-D)\bigr),
                         \qquad D=1,2,\ldots.
\]
Each has rank four, signature $(3,1)$, and a primitive
vector of positive square one. Its discriminant has
absolute value $D$. An integral isometry preserves
the determinant, because a change of integral basis
has determinant $\pm1$. Hence these lattices are
pairwise nonisometric.

They satisfy the rank and signature pattern of an
abstract Beauville--Bogomolov lattice, but they are
not asserted to occur as the second cohomology of
holomorphic symplectic manifolds. In particular no
geometric realizability or Fujiki constant is
attached to them. Their role is to disprove the
purely arithmetic implication that bounded
$b_2$, even with a small positive vector,
forces finitely many possible lattices.

For a more restrictive arithmetic sufficient
condition, a bound on all entries of a Gram
matrix in some integral basis, together with
a rank bound, does give finitely many lattices:
there are only finitely many such integer
matrices. However, choosing such bases with
uniform entry bounds is substantial information.
Changing to a real orthonormal basis loses
the integral lattice and cannot supply it.

\paragraph{A complete known rank bound in complex dimension four.}
For an irreducible holomorphic symplectic fourfold,
Salamon's relation and the injectivity of
$\operatorname{Sym}^2H^2(X,\mathbb Q)\to H^4(X,\mathbb Q)$
give, respectively,
\[
 b_4=10b_2+46-b_3,\qquad
              \frac{b_2(b_2+1)}2\leq b_4.
\]
These are established geometric inputs, recorded in
the discussion of Guan's bound [E1].
Combining them gives
\[
 b_2^2-19b_2-92+2b_3\leq0.
\]
Since $b_3\geq0$ and
$b_2^2-19b_2-92=(b_2-23)(b_2+4)$,
one obtains $b_2\leq23$. At $b_2=23$
the same inequality forces $b_3=0$ and
$b_4=276$. The standard Hilbert-square
example realizes these numerical values.

This is a genuinely useful restriction in one
dimension, not a finiteness proof in that
dimension. The preceding lattice family shows
the logical reason: rank control alone does
not control discriminants, intersection data
or deformation types. Nor may the fourfold
cohomological identity be reused unchanged
in higher complex dimensions.

\paragraph{Deformation invariants versus varying polarizations.}
A smooth proper family is locally differentiably
trivial, so its integral cohomology ring and
characteristic numbers are locally constant.
This explains why differing suitable topological
invariants can obstruct deformation equivalence.
It does not make their possible values bounded
when only the complex dimension is fixed.

Conversely unbounded polarization degree is not
evidence of infinitely many unpolarized types.
On one fixed projective manifold with ample
bundle $A$, all powers $A^k$ are ample and
\[
                         (A^k)^{2n}=k^{2n}A^{2n}.
\]
These degrees tend to infinity while the
underlying manifold has not changed at all.
The sufficient hypothesis
\eqref{a421:degree} only asks for the
existence of one controlled polarization
somewhere in each deformation type, not
a bound for every polarization.

The same distinction matters for families:
a parameter space of positive dimension
may describe many complex structures
inside one deformation class. An infinite
set of points in a moduli space is not
an infinite set of connected deformation
types. A negative answer would need an
infinite sequence in one fixed dimension
with an actual invariant separating the
deformation classes.

\paragraph{Why extra fibration data cannot be supplied for free.}
If a Lagrangian fibration is already available,
the isotropic class and mixed intersections
can make boundedness arguments much more
concrete. But the catalogue assumes neither
a fibration nor projectivity. Passing to
a deformation where an isotropic vector
has Hodge type $(1,1)$ is not the same as
showing that its line bundle is semiample.
The relevant semiampleness assertion is
itself a separate input discussed in [S1].
Even granting it would leave uniform
control of the invariants required by the
boundedness theorem.

Thus the calculation \eqref{a421:mixed}
is a consistency check and a conditional
tool. It must not be promoted to an
existence theorem for a map to a base.

\paragraph{Verification, first gap, and outcome.}
The degree and mixed coefficients were checked
for small $n$, including the factor
$\binom{2n}{n}$. The lattice determinant,
signature and primitive positive vector
were checked independently, and the
fourfold polynomial factors exactly as
displayed. None of these finite checks
asserts geometric realization of an
abstract lattice.

The first general gap is a uniform
boundedness input such as
\eqref{a421:degree}, or another mechanism
giving only finitely many admissible
topological and geometric possibilities
in each dimension.
\textbf{UNRESOLVED.} The conditional
polarization reduction, arithmetic
countertest and fourfold rank deduction
are established. They do not prove
unrestricted finiteness of deformation
types or produce an infinite family
of distinct types.

\subsection{Sources}
\begin{itemize}
\item[S1]Kamenova: Survey of finiteness results for hyperkähler manifolds.
\url{https://arxiv.org/pdf/1607.03215}.
\textit{Formulation source.}
\item[E1]Daniel Guan, \emph{On the Betti numbers of irreducible compact hyperk\"ahler manifolds of complex dimension four}, Mathematical Research Letters 8 (2001), 663--669.
\url{https://intlpress.com/api/bgcloud-front/resource/pdf/volume/1806606462270992385-1806606462270992385-66352cac07b0323617eb478779155915.pdf}.
\textit{Primary source for the known fourfold Betti restrictions and their cohomological inputs; consulted 27 September 2026. No fixed-dimension deformation classification is inferred from the rank bound.}
\end{itemize}
\end{document}
